\documentclass[12pt]{article}
\usepackage[a4paper,margin=1in]{geometry}
\usepackage[utf8]{inputenc}
\usepackage[T1]{fontenc}
\usepackage{amsmath,amssymb}
\usepackage{textcomp}
\usepackage{parskip}
\begin{document}
\begin{center}
{\Large \textbf{2003 STEP I}}
\end{center}

\subsection*{Page 1}
\begin{flushleft}
STEP I, 2003 2\par
\bigskip\textbf{Section A: Pure Mathematics}\par
n\par
1 It is given that r 2 can be written in the form pn 3 + qn 2 + rn + s , where p , q , r and s\par
\[r=\]
P\par
-1\par
\[are numbers. By setting n = 1, 0, 1 and 2, obtain four equations that must be satisfied by\]
-\par
p , q , r and s and hence show that\par
n\par
2 1\par
\[r = n(n + 1)(2n + 1) .\]
\[rP =0\]
n\par
3 4 3 2\par
Given that r can be written in the form an + bn + cn + dn + e , show similarly that\par
\[r=\]
P\par
-2\par
n\par
3 1 2 2\par
\[r = n (n + 1) .\]
\[rP =0\]
\bigskip\textbf{Question 2.} The first question on an examination paper is:\par
1 1 1\par
\[Solve for x the equation = + .\]
x a b\par
where (in the question) a and b are given non-zero real numbers. One candidate writes x = a+b\par
as the solution. Show that there are no values of a and b for which this will give the correct\par
answer.\par
The next question on the examination paper is:\par
1 1 1 1\par
\[Solve for x the equation = + + .\]
x a b c\par
where (in the question) a , b and c are given non-zero numbers. The candidate uses the same\par
technique, giving the answer as x = a + b + c . Show that the candidate's answer will be correct\par
if and only if a , b and c satisfy at least one of the equations a + b = 0 , b + c = 0 or c + a = 0 .\par
\bigskip\textbf{Question 3.} (i) Show that 2 sin( 1 θ) = sin θ if and only if sin( 1 θ) = 0 .\par
2 2\par
\textbf{(ii)} Solve the equation 2 tan( θ) = tan θ .\par
\textbf{(iii)} Show that 2 cos( θ) = cos θ if and only if θ = (4n + 2)\pi 2φ where φ is defined by\par
\[cos φ = 1 (\sqrt\{3\} 1) , 0 6 φ 6 \pi/2, and n is any integer. ±\]
-\par
\end{flushleft}
\subsection*{Page 2}
\begin{flushleft}
STEP I, 2003 3\par
\bigskip\textbf{Question 4.} Solve the inequality\par
sin θ + 1\par
6 1\par
cos θ\par
\[where 0 6 θ < 2\pi and cos θ = 0 .\]
\[5 (i) In the binomial expansion of (2x + 1/x 2 ) 6 for x = 0, show that the term which is\]
independent of x is 240.\par
Find the term which is independent of x in the binomial expansion of (ax\par
+ b/x\par
)\par
5n\par
.\par
\textbf{(ii)} Let f (x) = (x\par
+ 3x\par
)\par
\[1/2\]
. By considering the expansion of (1 + 3/x)\par
\[1/2\]
show that the\par
\[term which is independent of x in the expansion of f (x) in powers of 1/x , for x > 3 ,\]
| |\par
\[is 27/16 .\]
Show that there is no term independent of x in the expansion of f (x) in powers of x ,\par
\[for x < 3 .\]
| |\par
6 Evaluate the following integrals, in the different cases that arise according to the value of the\par
positive constant a :\par
\textbf{(i)} dx ;\par
\[Z^0 x^2 + (a + 2)x + 2a\]
\textbf{(ii)} du .\par
\[Z^1 u^2 + au + a 1\]
-\par
\bigskip\textbf{Question 7.} Let k be an integer satisfying 0 6 k 6 9 . Show that 0 6 10k k 2 6 25 and that k(k 1)(k + 1)\par
- -\par
is divisible by 3 .\par
For each 3-digit number N , where N > 100, let S be the sum of the hundreds digit, the square\par
\[of the tens digit and the cube of the units digit. Find the numbers N such that S = N .\]
\[[Hint: write N = 100a + 10b + c where a , b and c are the digits of N .]\]
\end{flushleft}
\subsection*{Page 3}
\begin{flushleft}
STEP I, 2003 4\par
\bigskip\textbf{Question 8.} A liquid of fixed volume V is made up of two chemicals A and B . A reaction takes place in\par
which A converts to B . The volume of A at time t is xV and the volume of B at time t is\par
\[yV where x and y depend on t and x + y = 1 . The rate at which A converts into B is given\]
by kV xy , where k is a positive constant. Show that if both x and y are strictly positive at\par
the start, then at time t\par
D ekt\par
\[y = ,\]
1 + D ekt\par
where D is a constant.\par
Does A ever completely convert to B ? Justify your answer.\par
\end{flushleft}
\subsection*{Page 4}
\begin{flushleft}
STEP I, 2003 5\par
\bigskip\textbf{Section B: Mechanics}\par
\bigskip\textbf{Question 9.} A particle is projected with speed V at an angle θ above the horizontal. The particle passes\par
through the point P which is a horizontal distance d and a vertical distance h from the point\par
of projection. Show that\par
2kh\par
\[T 2kT + + 1 = 0 ,\]
- d\par
V\par
\[where T = tan θ and k = .\]
gd\par
2 2\par
Show that, if kd > h + \sqrt(h + d ) , there are two distinct possible angles of projection.\par
\[Let these two angles be α and β. Show that α + β = \pi arctan(d/h) .\]
-\par
10 ABCD is a uniform rectangular lamina and X is a point on BC . The lengths of AD, AB\par
and BX are p , q and r respectively. The triangle ABX is cut off the lamina. Let (a, b) be the\par
position of the centre of gravity of the lamina, where the axes are such that the coordinates\par
of A , D and C are (0, 0) , (p, 0) and (p, q) respectively. Derive equations for a and b in terms\par
of p , q and r .\par
When the resulting trapezium is freely suspended from the point A , the side AD is inclined\par
at 45 \textasciicircum{}\circ below the horizontal. Show that r = q q2 3pq + 3p2 . You should justify carefully\par
- -\par
the choice of sign in front of the square root. p\par
\bigskip\textbf{Question 11.} A smooth plane is inclined at an angle α to the horizontal. A and B are two points a distance d\par
apart on a line of greatest slope of the plane, with B higher than A. A particle is projected\par
up the plane from A towards B with initial speed u, and simultaneously another particle is\par
released from rest at B . Show that they collide after a time d/u .\par
The coefficient of restitution between the two particles is e and both particles have mass m .\par
1 2 2\par
Show that the loss of kinetic energy in the collision is mu 1 e .\par
-\par
\end{flushleft}
\subsection*{Page 5}
\begin{flushleft}
STEP I, 2003 6\par
\bigskip\textbf{Section C: Probability and Statistics}\par
\bigskip\textbf{Question 12.} In a bag are n balls numbered 1, 2, . . . , n . When a ball is taken out of the bag, each ball is\par
equally likely to be taken.\par
\textbf{(i)} A ball is taken out of the bag. The number on the ball is noted and the ball is\par
replaced in the bag. The process is repeated once. Explain why the expected value\par
of the product of the numbers on the two balls is\par
n n\par
rs\par
\[n^2\]
\[Xr=1 Xs=1\]
and simplify this expression.\par
\textbf{(ii)} A ball is taken out of the bag. The number on the ball is noted and the ball is not\par
replaced in the bag. Another ball is taken out of the bag and the number on this ball\par
is noted. Show that the expected value of the product of the two numbers is\par
(n + 1)(3n + 2)\par
.\par
n n\par
1 1\par
\[Note: r = n(n + 1) and r 2 = n(n + 1)(2n + 1) .\]
2 6\par
\[Xr=1 Xr=1\]
13 If a football match ends in a draw, there may be a "penalty shoot-out". Initially the teams\par
each take 5 shots at goal. If one team scores more times than the other, then that team wins.\par
If the scores are level, the teams take shots alternately until one team scores and the other\par
team does not score, both teams having taken the same number of shots. The team that\par
scores wins.\par
Two teams, Team A and Team B, take part in a penalty shoot-out. Their probabilities of\par
scoring when they take a single shot are p and p respectively. Explain why the probability\par
A B\par
α of neither side having won at the end of the initial 10-shot period is given by\par
5 2\par
\[α = (1 p )i(1 p )i p 5-ip 5-i.\]
\lambdai\mu - A - B A B\par
\[Xi=0\]
2α\par
\[Show that the expected number of shots taken is 10 + , where β = p + p 2p p .\]
A B A B\par
β -\par
\end{flushleft}
\subsection*{Page 6}
\begin{flushleft}
STEP I, 2003 7\par
14 Jane goes out with any of her friends who call, except that she never goes out with more than\par
two friends in a day. The number of her friends who call on a given day follows a Poisson\par
distribution with parameter 2. Show that the average number of friends she sees in a day\par
-2\par
is 2 4 e .\par
-\par
Now Jane has a new friend who calls on any given day with probability p. Her old friends call\par
as before, independently of the new friend. She never goes out with more than two friends in\par
a day. Find the average number of friends she now sees in a day.\par
\end{flushleft}
\end{document}
